\documentclass[11pt]{article}
\usepackage[utf8]{inputenc}
\usepackage[T1]{fontenc}
\usepackage[margin=1in]{geometry}
\usepackage{amsmath}
\usepackage{booktabs}
\usepackage{graphicx}
\usepackage{hyperref}
\usepackage{natbib}
\usepackage{xcolor}
\usepackage{listings}
\usepackage{lmodern}
\usepackage{microtype}

\lstset{
  basicstyle=\ttfamily\small,
  columns=fullflexible,
  breaklines=true,
  frame=single,
  backgroundcolor=\color{black!3}
}

\title{leaklens: A Unified Toolkit for Benchmark Contamination Detection}
\author{Shreyas K Chandrahas}
\date{}

\begin{document}
\maketitle

\begin{abstract}
Benchmark contamination detection methods are individually well studied --
n-gram overlap against pretraining corpora, Min-K\% Prob membership
inference \citep{shi2023detecting}, the zlib-entropy perplexity baseline
\citep{carlini2021extracting}, guided completion, and order-sensitivity
probes in the tradition of the GPT-3 contamination appendix
\citep{brown2020language} -- but no maintained, standard implementation
packages them behind one interface, so every paper that wants to check for
contamination either reimplements a detector or waves hands. We present
\texttt{leaklens}, a plugin-based toolkit implementing six such detectors:
n-gram overlap against the public infini-gram index
\citep{liu2024infinigram}, guided completion, order canary, Min-K\% Prob,
a perplexity-vs-compressibility gap, and paraphrase gap. Detectors are
explicit about their own applicability: logprob-based methods (Min-K\%
Prob, perplexity gap, paraphrase gap) require a local model and report
\texttt{applicable=False} rather than silently failing against API-only
judges, and every report card states which detectors ran, which were
skipped, and why. We validate the toolkit with a real-hardware calibration
pilot: two LoRA adapters fine-tuned on disjoint 15-item subsets of GSM8K
(one ``contaminated'', one ``clean'') using a local Qwen2.5-0.5B model,
then scored with every logprob- and completion-based detector on the
contaminated subset. Guided completion, Min-K\% Prob, and the
perplexity-compressibility gap all separate the fine-tuned-on model from
the not-fine-tuned-on model cleanly (AUC 0.92, 1.00, and 1.00
respectively, on this small $n=15$ pilot); a fourth detector,
order-sensitivity, scored at chance, and we trace this to a genuine
methodological artifact in the calibration's own training-data design
(documented in Section~\ref{sec:calibration}) rather than to a flaw in the
detection method itself. All six of the toolkit's dataset benchmark
adapters (MMLU, GSM8K, HumanEval, ARC, HellaSwag, TruthfulQA) are verified
against the live HuggingFace \texttt{datasets-server} API rather than
assumed from memory. \texttt{leaklens} is available at
\url{https://github.com/Shreyaskc/leaklens}.
\end{abstract}

\section{Introduction}

Benchmark contamination -- a model having seen a benchmark's items, or
close paraphrases of them, during training -- is one of the central
validity threats in LLM evaluation. A growing body of work has produced
individual detection methods: substring/n-gram overlap against known
pretraining corpora, membership-inference statistics computed from a
model's own token log-probabilities \citep{shi2023detecting}, extraction
and compressibility-based probes \citep{carlini2021extracting}, and
behavioral probes that ask whether a model can complete a benchmark item
verbatim or recover its canonical ordering \citep{brown2020language}. Each
method has a paper; none has a maintained, general-purpose implementation
that a researcher auditing a new benchmark--model pair can simply install
and run. The result is that contamination checks, when they appear in
papers at all, are usually a single ad hoc n-gram search rather than the
fuller battery the literature already supports.

\texttt{leaklens} packages six such methods behind one API,
\texttt{leaklens.audit(model, benchmark)}, which returns a report card
listing, per detector, whether it ran, its aggregate score, and (for
detectors that were not applicable to the given model) why. The library
is deliberately an aggregator, not a competitor to the underlying methods:
every detector's docstring and every report card cites the paper the
method comes from, so leaklens is cited alongside the primitives it
wraps, not instead of them.

\section{The leaklens Toolkit}

\subsection{Plugin interface}

Every detector implements a common \texttt{Detector} interface with a
\texttt{requires\_logprobs} flag and a \texttt{run(model, benchmark)}
method. A \texttt{ModelInterface} abstraction exposes \texttt{generate()}
(supported by any model, local or API) and \texttt{token\_logprobs()}
(supported only by local models with weight access). Before running a
detector, \texttt{leaklens.audit()} checks
\texttt{detector.applicable(model)}; if a logprob-only detector is asked
to run against an API-only model, the report card records
\texttt{applicable=False} with an explicit reason rather than crashing or
silently omitting the row. This matters for reproducibility: a report card
should show what was \emph{not} checked, not just what was.

\subsection{Detectors}

\textbf{N-gram overlap.} Queries the public infini-gram API
\citep{liu2024infinigram} for exact-match counts of a benchmark item's
text in a named pretraining-corpus index (Dolma, the Pile, C4, or
RedPajama). This detector is deliberately model-independent -- it answers
whether an item's text is present in a \emph{public corpus}, not whether
the specific model under audit was trained on that exact corpus, and its
report-card entry states this scope explicitly. The endpoint and response
schema were verified against the live API during development rather than
assumed.

\textbf{Guided completion.} Prompts the model with a prefix of a
benchmark item (a configurable fraction of its length) and scores the
completion's similarity to the true continuation via a normalized
sequence-matching ratio. High-fidelity verbatim reproduction of a
continuation the model was never shown at inference time is a
memorization signal in the tradition of training-data-extraction attacks
\citep{carlini2021extracting}.

\textbf{Order canary.} For each adjacent pair of items in a benchmark's
canonical published order, prompts the model with the tail of item $i$
and scores whether the completion matches the head of item $i{+}1$ (the
real next item) versus a randomly reassigned ``fake next'' item (a
shuffled control). A model that predicts the real next item reliably
better than the shuffled control has memorized the benchmark's exact
sequence, independent of item content -- a content-independent
contamination signal in the tradition of GPT-3's own contamination
appendix \citep{brown2020language}.

\textbf{Min-K\% Prob.} Following \citet{shi2023detecting}, computes the
mean log-probability of the $k$\% least-likely tokens in an item's text.
Text a model has memorized tends to have a less-negative worst-case
token probability than genuinely unseen text, since exact repetition
raises even the model's least confident predictions.

\textbf{Perplexity-compressibility gap.} Following the zlib-entropy
baseline from \citet{carlini2021extracting}, divides a model's
cross-entropy on an item by the item's zlib-compressed byte length. Raw
perplexity alone confounds memorization with an item simply being
low-entropy; dividing by a model-independent compressibility measure
controls for that confound.

\textbf{Paraphrase gap.} Compares a model's cross-entropy on an item's
original text versus a semantically-equivalent paraphrase. A model that
finds the exact original phrasing much easier than an equivalent
paraphrase is showing surface-form memorization rather than task
competence. This detector's logic is implemented and tested, but
leaklens does not yet ship pre-generated, human-checked paraphrases for
its six built-in benchmarks -- a documented data-construction gap, not a
silent omission; callers may supply their own paraphrases via a field on
each benchmark item.

\subsection{Benchmark adapters}

Built-in adapters wrap six widely used benchmarks -- MMLU, GSM8K,
HumanEval, ARC, HellaSwag, and TruthfulQA -- as sources of items with a
fixed, canonical ordering (required for the order-canary detector to be
meaningful). Every adapter's HuggingFace dataset repository ID, config
name, split name, and column names were checked directly against the live
\texttt{datasets-server} API during development, not assumed from prior
familiarity with these datasets; the checks are re-run as part of the
test suite so upstream schema drift is caught rather than silently
producing empty or malformed items.

\section{Calibration Pilot}
\label{sec:calibration}

The scientific claim a contamination detector makes -- ``this score
indicates the model saw this item during training'' -- is only as good as
its calibration against known ground truth. We ran a small, fast
calibration pilot entirely on local hardware (an Apple M4 Pro, no rented
or cloud GPU): two LoRA adapters were fine-tuned from the same base model,
\texttt{Qwen2.5-0.5B-Instruct} (4-bit), on two disjoint 15-item subsets of
GSM8K test questions -- one subset (``contaminated'') the model was
repeatedly fine-tuned on, one (``clean'') it never saw. Each of the four
model-dependent detectors was then run against the contaminated subset's
15 items using both fine-tuned models, and we computed the AUC-ROC of
each detector's score at distinguishing the contaminated (member) model
from the clean (non-member) model on the same items.

\begin{table}[t]
\centering
\begin{tabular}{lc}
\toprule
Detector & AUC \\
\midrule
Guided completion & 0.92 \\
Min-K\% Prob & 1.00 \\
Perplexity--compressibility gap & 1.00 \\
Order canary & 0.50 \\
\bottomrule
\end{tabular}
\caption{Calibration pilot AUC-ROC, $n=15$ items, one base model, one
fine-tune pass per condition. Ngram overlap and paraphrase gap are not
included -- see Section~\ref{sec:scope}.}
\label{tab:calibration}
\end{table}

Table~\ref{tab:calibration} reports the results. Three detectors --
guided completion, Min-K\% Prob, and the perplexity-compressibility gap
-- cleanly separate the two models on this pilot. The fourth,
order canary, scored at chance (AUC $\approx 0.50$), with both models'
completions at the relevant prompt boundary being an empty string. We
traced this to a real artifact of this calibration's own training-data
construction, not a flaw in the underlying method: each individual GSM8K
item was included in the fine-tuning data as its own training example
(implicitly teaching the model ``predict end-of-sequence after this exact
tail''), while adjacent-item pairs were \emph{also} included as
continuous training examples (teaching ``continue into the next item'')
at the exact same text boundary. These two training signals conflict at
inference time with no way for the model to disambiguate which one
applies, and it appears to default to stopping. An isolated calibration
design -- training only on the continuous concatenated sequence, without
the individual-item examples -- would be needed to calibrate order canary
cleanly, and is left for a follow-up run. We report this null result and
its diagnosis rather than omit the detector or present only the three
detectors that happened to work.

\subsection{Scope of the calibration}
\label{sec:scope}

Two of the six detectors are not covered by this calibration, for
principled reasons rather than oversight. N-gram overlap queries a fixed,
public pretraining-corpus index and is architecturally independent of any
model under audit (see Section 2.2); a local LoRA fine-tune on GSM8K items
has no bearing on whether those items independently appear in a
large public web corpus, so this pilot's fine-tuning design cannot
calibrate it. Paraphrase gap requires paraphrase data this release does
not yet ship (Section 2.2).

We also note plainly that this is a single small pilot -- one base
model, $n=15$ items, one fine-tuning run per condition -- reported as an
indicative first pass, not the larger multi-model calibration suite (with
true/false positive rates across many models and content types) that a
mature calibration study would require.

\section{Validation}

The package includes 74 unit and integration tests (99\% statement
coverage), including tests that exercise the real, live infini-gram API
response schema (not a hand-written mock of an assumed schema), the real
MLX-based logprob computation validated against expected model behavior
(a repeated sentence scores a substantially higher log-probability on its
second occurrence than its first; character-level gibberish scores
reliably lower than natural language), and the calibration pipeline
described in Section~\ref{sec:calibration}, which was itself run twice:
a first run produced NaN training loss partway through fine-tuning,
traced to a single training example roughly an order of magnitude longer
than the rest of its batch destabilizing gradients; the fix (splitting
that one long example into shorter adjacent-pair examples) is reflected
in the released calibration script.

\section{Limitations}

The calibration pilot in Section~\ref{sec:calibration} is small-scale by
design (one base model, $n=15$ items) and should be read as a sanity
check that leaklens's detectors behave in the expected direction, not as
a definitive true/false-positive-rate characterization suitable for
setting a universal decision threshold. The boilerplate-phrase and
artifact-signature detection this toolkit's broader design contemplates
is not yet part of the released detector set. Paraphrase gap's data gap
(Section 2.2) means it is currently usable only by callers who supply
their own paraphrases. Order canary's null result in this pilot
(Section~\ref{sec:calibration}) means we do not yet have a clean
calibration data point for that detector; the diagnosis suggests a fix,
but the fix itself has not yet been run.

\section{Ethics Statement}

This work is a measurement toolkit; it does not collect, redistribute, or
generate new human-subject data. The infini-gram queries send benchmark
item text (already public benchmark content) to a third-party public API;
no private or user data is involved. The calibration pilot fine-tunes a
small open-weight model on public GSM8K questions already released under
an open license, entirely on local hardware.

\section{Availability}

\texttt{leaklens} is available at
\url{https://github.com/Shreyaskc/leaklens} under the MIT license.

\section{Conclusion}

We presented \texttt{leaklens}, a plugin-based contamination-detection
toolkit unifying six methods from the literature behind one interface,
with explicit per-detector applicability reporting and six verified
benchmark adapters. A local-hardware calibration pilot confirms that
three of four model-dependent detectors cleanly separate a
fine-tuned-on item set from a held-out one on a small sample, and traces
the fourth detector's null result to a specific, fixable artifact of the
calibration's own design rather than papering over it.

\bibliographystyle{plainnat}
\bibliography{references}

\end{document}
